\documentclass[12pt]{article}
\usepackage{amsmath,amssymb,amsfonts}
\usepackage[margin=1in]{geometry}

\begin{document}

\begin{center}
{\Large\bfseries Chapter 10, Problem 16}
\end{center}

\noindent\textbf{Problem.} Find the characters and irreducible representations of the group discussed in Problem 10.3. \\
\textit{[Hint: For the two-dimensional representation, the result of Problem 10.15 may be useful.]}

\bigskip
\noindent\textbf{Solution.}

The group from Problem 3 is $D_4$ (the dihedral group of order 8). Its 8 elements fall into 5 conjugacy classes:
\begin{align*}
&\mathcal{C}_1 = \{e\}, \\
&\mathcal{C}_2 = \{r, r^3\} \quad (\text{rotations by } 90^\circ, 270^\circ), \\
&\mathcal{C}_3 = \{r^2\} \quad (\text{rotation by } 180^\circ), \\
&\mathcal{C}_4 = \{\sigma_1, \sigma_2\} \quad (\text{reflections through edge midpoints}), \\
&\mathcal{C}_5 = \{\sigma_d, \sigma_{d'}\} \quad (\text{reflections through diagonals}).
\end{align*}

With 5 classes, there are 5 irreducible representations. Since $\sum n_i^2 = 8$ and $n_i \geq 1$, the only solution is $1^2 + 1^2 + 1^2 + 1^2 + 2^2 = 8$. So there are four 1-dimensional and one 2-dimensional irreducible representations.

\medskip
\noindent\textbf{One-dimensional representations:}

Each 1D representation is a homomorphism $D_4 \to \{1, -1\}$ (since elements have finite order). Since $r^4 = e$, we need $D(r)^4 = 1$, and for 1D real representations $D(r) = \pm 1$.

If $D(r) = 1$: then $D(r^2) = 1$. The reflections satisfy $\sigma^2 = e$, so $D(\sigma) = \pm 1$.

If $D(r) = -1$: then $D(r^2) = 1$, $D(r^3) = -1$. Again $D(\sigma) = \pm 1$.

The constraint from the group structure $r\sigma_1 = \sigma_d r$ (i.e., the reflections are related by rotations) gives $D(r)D(\sigma_1) = D(\sigma_d)D(r)$, so if $D(r) = 1$ then $D(\sigma_1) = D(\sigma_d)$ is unconstrained; but the relation $r\sigma_1 r^{-1} = \sigma_2$ means $D(\sigma_1) = D(\sigma_2)$ and $r\sigma_d r^{-1} = \sigma_{d'}$ means $D(\sigma_d) = D(\sigma_{d'})$.

The four 1D representations are:

\[
\begin{array}{c|ccccc}
 & \mathcal{C}_1 & \mathcal{C}_2 & \mathcal{C}_3 & \mathcal{C}_4 & \mathcal{C}_5 \\
 & \{e\} & \{r,r^3\} & \{r^2\} & \{\sigma_1,\sigma_2\} & \{\sigma_d,\sigma_{d'}\} \\
\hline
\Gamma_1 & 1 & 1 & 1 & 1 & 1 \\
\Gamma_2 & 1 & 1 & 1 & -1 & -1 \\
\Gamma_3 & 1 & -1 & 1 & 1 & -1 \\
\Gamma_4 & 1 & -1 & 1 & -1 & 1 \\
\end{array}
\]

\medskip
\noindent\textbf{Two-dimensional representation:}

Using orthogonality of characters to find $\chi^{(5)}$. Let $\chi^{(5)} = (2, a, b, c, d)$ for classes $\mathcal{C}_1, \ldots, \mathcal{C}_5$.

Orthogonality with $\Gamma_1$: $1\cdot 2 + 2\cdot a + 1\cdot b + 2\cdot c + 2\cdot d = 0 \implies 2a + b + 2c + 2d = -2$.

Orthogonality with $\Gamma_2$: $2 + 2a + b - 2c - 2d = 0 \implies 2a + b = 2c + 2d - 2$.

From these two: $4c + 4d = 0$, so $d = -c$. Then $2a + b = -2$.

Orthogonality with $\Gamma_3$: $2 - 2a + b + 2c - 2d = 0 \implies -2a + b + 4c = -2$.

Combined with $2a + b = -2$: $-4a + 4c = 0$, so $c = a$, hence $d = -a$.

From $2a + b = -2$: $b = -2 - 2a$.

Normalization: $1\cdot 4 + 2a^2 + (2+2a)^2 + 2a^2 + 2a^2 = 8$.
\[
4 + 2a^2 + 4 + 8a + 4a^2 + 2a^2 + 2a^2 = 8 \implies 10a^2 + 8a = 0 \implies a(10a + 8) = 0.
\]
So $a = 0$ or $a = -4/5$. Since the dimension is 2 and $|\chi(g)| \leq 2$, we need $|b| = |{-2-2a}| \leq 2$. For $a = 0$: $b = -2$, $c = 0$, $d = 0$. Check: $|-2| = 2 \leq 2$. \checkmark

The character table is:
\[
\begin{array}{c|ccccc}
 & \{e\} & \{r,r^3\} & \{r^2\} & \{\sigma_1,\sigma_2\} & \{\sigma_d,\sigma_{d'}\} \\
 & 1 & 2 & 1 & 2 & 2 \\
\hline
\Gamma_1 & 1 & 1 & 1 & 1 & 1 \\
\Gamma_2 & 1 & 1 & 1 & -1 & -1 \\
\Gamma_3 & 1 & -1 & 1 & 1 & -1 \\
\Gamma_4 & 1 & -1 & 1 & -1 & 1 \\
\Gamma_5 & 2 & 0 & -2 & 0 & 0 \\
\end{array}
\]

\medskip
\noindent\textbf{Explicit 2D representation matrices:}

The natural 2D representation comes from the action of $D_4$ on the plane. With $r$ = rotation by $90^\circ$ and $\sigma_1$ = reflection through the $x$-axis:

\[
D(r) = \begin{pmatrix} 0 & -1 \\ 1 & 0 \end{pmatrix}, \quad
D(r^2) = \begin{pmatrix} -1 & 0 \\ 0 & -1 \end{pmatrix}, \quad
D(r^3) = \begin{pmatrix} 0 & 1 \\ -1 & 0 \end{pmatrix},
\]
\[
D(\sigma_1) = \begin{pmatrix} 1 & 0 \\ 0 & -1 \end{pmatrix}, \quad
D(\sigma_2) = \begin{pmatrix} -1 & 0 \\ 0 & 1 \end{pmatrix},
\]
\[
D(\sigma_d) = \begin{pmatrix} 0 & 1 \\ 1 & 0 \end{pmatrix}, \quad
D(\sigma_{d'}) = \begin{pmatrix} 0 & -1 \\ -1 & 0 \end{pmatrix}.
\]

One can verify these traces match the character table above.

\end{document}
